\documentclass[../../../thesis.tex]{subfiles}
\begin{document}
\subsection{Jazyk a gramatika}
Záverečná práca na FEI STU v~Bratislave musí byť napísaná
buď po slovensky alebo po anglicky.
Ak je jazyk práce angličtina, musí po závere nasledovať
rezumé v~slovenskom jazyku.

Záverečná práca univerzitného štúdia sa vyznačuje
vysokou jazykovou úrovňou.
Gramatické a štylistické chyby sú neprípustné.
Študent by mal tejto stránke diela venovať patričnú
pozornosť a~podľa možností nechať rukopis prejsť
kvalifikovanou jazykovou kontrolou.
Najmä bakalárska práca predstavuje v~živote väčšiny študentov
prvý rozsiahlejší autorský útvar,
ktorý má významný vplyv na jeho ďalší život a~kariéru.

Aj keď väčšina textových editorov dokáže odhaľovať preklepy,
neporadí si s~komplikovanejšou gramatikou a~štylistikou.
Treba sa riadiť najmä pravidlami slovenského pravopisu,
slovníkmi slovenského jazyka a~ďalšími zdrojmi,
ktoré možno nájsť na webových stránkach
Jazykovedného ústavu Ľudovíta Štúra SAV.%
\footnote{\href{https://www.juls.savba.sk/}{\tt www.juls.savba.sk}}
Využiť môžeme aj jazykovú poradňu,
ktorú poskytuje ústav bezplatne a~to buď telefonicky alebo
prostredníctvom emailovej komunikácie.
Cenným zdrojom informácií môže byť aj Jazyková poradňa
denníka SME v spolupráci
s~Jazykovedným ústavom Ľudovíta Štúra SAV%
\footnote{\href{https://jazykovaporadna.sme.sk/}{\tt jazykovaporadna.sme.sk}}
alebo online slovníky slovenského jazyka,%
\footnote{\href{https://slovnik.juls.savba.sk/}{\tt slovnik.juls.savba.sk}}
prípadne národný jazykový korpus.%
\footnote{\href{https://korpus.sk/}{\tt korpus.sk}}

Pri písaní práce dbáme najmä na pravopisné javy ako sú písanie
tvrdého a~mäkkého y/i vo vybraných slovách,
v~príponách a~koncovkách pri skloňovaní
(pekný muž, ale pekní muži),
v~číslovkách (rozprávali sme sa so siedmimi v~poradí
-- skončili siedmi v~poradí,
ale hrali sme sa so siedmymi deťmi -- detí bolo sedem), atď.
Rovnako dôležité je správne písanie rodov,
skloňovanie a časovanie.

Veľmi komplexná a~dôležitá zložka gramatiky
je písanie čiarok v~súvetiach.

Popri gramatike je podstatná aj štylistická tvorba viet,
ktorú musí študent univerzity zvládať na vysokej úrovni.

\subsubsection{Delenie slov}
Tzv. \emph{textové procesory} ako MS Word, LibreOffice a~Apache OpenOffice
ponúkajú automatické delenie slov na konci riadka.
Systém na sadzbu textu \LaTeX\ má túto funkciu automaticky zapnutú a~jej slovenská lokalizácia je veľmi kvalitne spracovaná.

Vo veľkej väčšine prípadov je delenie
v~súlade s~pravidlami jazyka.
Môžu sa vyskytnúť sporné okolnosti,
kedy počítač nerozdelí slovo správne.
Väčšinou máme možnosť do procesu zasiahnuť
a~ručne kontrolovať delenie slov na miestach,
s ktorými si softvér nevie poradiť.
Príkaz na preferované rozdelenie slova je \verb|\-|.
Napríklad slovo \verb|predstave\-nie| \LaTeX\
preferovane rozdelí v mieste prípony.

V každom prípade je žiadúce slová na konci riadka deliť
a\ túto možnosť nevypínať.
Prospieva to práci ako po technickej,
tak aj po estetickej stránke.
Odseky obsahujú menej dier,
textová oblasť stránky je vyplnená homogénnejšie,
čo prispieva k~lepšej čitateľnosti.
V prípade, že používame zarovnávanie do bloku tak,
ako aj v~tomto dokumente,
je prítomnosť dier v~odseku značne rušivá.
Ak používame zarovnanie textu doľava,
nepoužívanie delenia slov má vplyv na vznik tzv. riek,
čo je náhle striedanie dlhých a~krátkych riadkov.
Pravý okraj textu je nepekne zubatý.

Pravidlá rozdeľovania slov na konci riadka sú pomerne zložité.
Základné pravidlo, ktoré si pamätáme zo základnej školy,
je, že slová delíme na slabiky pred spoluhláskou alebo medzi
dvomi spoluhláskami.
Ak si nie sme istí, uprednostňujeme delenie v mieste,
kde sa ku koreňu slova pripájajú predpony alebo prípony,
prípadne v~mieste spojenia slov v~zloženom slove.

Pri slovách utvorených predponou alebo príponou
uprednostňujeme morfologické delenie
pred rozdelením koreňa slova.
Najskôr sa snažíme deliť slovo za predponou,
ak to nejde, skúsime to pred príponou.
Napríklad slovo \emph{predstavenie}
delíme na slabiky takto: \emph{pred-sta-ve-nie}.
Pri rozdeľovaní slov uprednostňujeme model
\emph{pred-stavenie}, výnimočne aj \emph{pred-stave-nie}.
V slove \emph{výklenok} sa uplatňuje pravidlo morfologického
delenia pred delením v~mieste zhluku spoluhlások.
Sylabická stavba tohto slova je \emph{vý-kle-nok},
nie \emph{výk-le-nok},
pretože slovo pozostáva z troch častí: predpony \emph{vý},
koreňa \emph{kle} a prípony \emph{nok}.
Mohli by sme namietať, že prípona je \emph{ok},
pomocou ktorej bolo vytvorené podstatné meno zo slovesa
klenúť alebo z~prídavného mena klenutý,
kde identifikujeme koreň \emph{klen}.
V skutočnosti je však príponou \emph{-nok}.
Morfológia je pomerne komplexná problematika,
a~nedokážeme tu obsiahnuť všetky jej detaily.
Väčšinou sa môžeme spoľahnúť na softvér,
že slová rozdelí správne.
V prípade pochybností využijeme externé pomôcky spomenuté
v~úvode tejto kapitoly.

Slová spojené spojovníkom rozdeľujeme v~mieste spojovníka tak,
že spojovník napíšeme na konci aj na začiatku riadka.
Slovo \emph{vedecko-pedagogický} môžeme rozdeliť takto:
\emph{ve-dec-ko-pe-da-go-gic-ký}.
Ak delenie padne na miesto spojenia slov,
rozdelíme ho nasledujúcim spôsobom:\\
\indent\emph{vedecko-}\\
\indent\emph{-pedagogický}

V šablóne rieši tento problém príkaz
\verb|\languageattribute{slovak}{split}|,
ktorý je súčasťou jazykového balíka \verb|babel|.

Nesprávne delenie slov sa v~práci zvyčajne objaví
len zriedkavo a~nemá vplyv na jej hodnotenie.
Netreba sa naň príliš sústrediť a~robiť si starosti.
Celkový vzhľad práce viac naruší vypnutie delenia slov,
než občasná malá chyba.

\subsubsection{Jednopísmenové predložky a spojky}
Hovoríme o predložkách k, o, v, s, z,
ktoré by nemali ostať osamotené na konci riadka.
Do tejto kategórie patria aj spojky a, i.
Jednopísmenové slová pripájame k~nasledujúcemu slovu pomocou
tzv. \emph{nedeliteľnej medzery},
čo je špeciálny netlačiteľný znak.
V kódovaní UTF-8 má číslo 00A0 (ASCII 160)
a~hovorí textovému procesoru,
že na tomto mieste nesmie byť za žiadnych okolností
koniec riadka.
V programe MS Word ho zadáme použitím klávesovej skratky
Ctrl-Shift-Medzera.
V \LaTeX-u\ zapíšeme nedeliteľnú medzeru s~premenlivou šírkou
ako symbol vlnovka (\verb|~|)
Napríklad slovné spojenie \emph{v~priestore} napíšeme takto:
\verb|v~priestore|.
Na rozdiel od Wordu, \LaTeX\ takéto medzery nevkladá automaticky
a treba to urobiť ručne.

Existuje viacero medzier, ktoré sú tiež nedeliteľné a~majú pevnú šírku.
Najpoužívanejšia tzv. úzka medzera a zapíšeme ju ako \verb|\,|.
Takýto typ medzery používame pri zápise hodnôt fyzikálnych veličín a vkladáme ju medzi číslo a jednotku.
\end{document}